\chapter{Variáveis: Guardando Informações}

\begin{objetivos}
\begin{itemize}
\item Entender o que é uma variável, como uma caixinha com nome;
\item Criar variáveis e mudar o valor delas;
\item Fazer um contador de pontos.
\end{itemize}
\end{objetivos}

\section{A caixinha com nome}

Imagine uma caixinha com uma etiqueta escrita \textbf{pontos}. Dentro dela cabe um número. Você pode olhar o que tem dentro, trocar por outro ou somar 1. Essa caixinha é uma \textbf{variável}.

\begin{vocabulario}[Vocabulário: variável]
\textbf{Variável} é um espaço na memória do computador com um \textbf{nome} e um \textbf{valor}. O valor pode \textbf{variar} (mudar) durante o jogo, e daí vem o nome!
\end{vocabulario}

\begin{pensaassim}[Pense assim: o placar do jogo]
O placar de um jogo de futebol começa em 0 a 0 e muda a cada gol. \textbf{Defina} é zerar o placar no começo da partida; \textbf{mude em 1} é marcar mais um gol. Se você esquecer de zerar no início, o placar da partida anterior continua lá!
\end{pensaassim}

\section{Criando uma variável}

Na categoria \bloco[scVariables]{Variáveis}, clique em \textbf{Criar uma variável}, dê um nome (por exemplo, \textbf{pontos}) e escolha \textbf{para todos os atores}. Pronto! A variável aparece no palco e ganha novos blocos.

\begin{center}
\renewcommand{\arraystretch}{1.8}
\begin{tabular}{@{}ll@{}}
\bloco[scVariables]{defina \menu{pontos} para \entrada{0}} & guarda um valor \\
\bloco[scVariables]{mude \menu{pontos} em \entrada{1}} & soma (ou subtrai, se for negativo) \\
\bloco[scVariables]{mostre a variável \menu{pontos}} & aparece no palco \\
\bloco[scVariables]{esconda a variável \menu{pontos}} & some do palco \\
\variavel{pontos} & o valor guardado \\
\end{tabular}
\end{center}

\begin{dica}
Dê nomes claros às variáveis: \textbf{pontos}, \textbf{vidas}, \textbf{tempo}. Evite \textbf{x1} ou \textbf{coisa}, pois daqui a uma semana você nem vai lembrar o que era!
\end{dica}

\section{Contando pontos}

\begin{missao}[Missão 13: O contador de cliques]
Crie a variável \textbf{pontos}. Cada vez que clicarem no gatinho, o placar sobe 1.
\end{missao}

\begin{codigo}[Placar de cliques]
\begin{pilha}
\bk[scEvents]{quando \bandeira\ for clicado}
\bk[scVariables]{defina \menu{pontos} para \entrada{0}}
\end{pilha}
\begin{pilha}
\bk[scEvents]{quando este ator for clicado}
\bk[scVariables]{mude \menu{pontos} em \entrada{1}}
\bk[scSound]{inicie o som \menu{Miau}}
\end{pilha}
\end{codigo}

\begin{cuidado}
Sempre \textbf{zere} a variável quando o jogo começar, com \bloco[scVariables]{defina \menu{pontos} para \entrada{0}}. Senão o valor do jogo anterior continua lá!
\end{cuidado}

\section{Variável + condição = jogo}

Variáveis ficam poderosas quando combinadas com decisões. Por exemplo, dar a vitória quando os pontos chegam a 10:

\begin{codigo}[Ganhou!]
\begin{pilha}
\bk[scEvents]{quando \bandeira\ for clicado}
\sbc[scControl]{sempre}
\sbc[scControl]{se \booleano{\variavel{pontos} = \entrada{10}} então}
\bk[scLooks]{diga \entrada{Você venceu!}}
\bk[scControl]{pare \menu{tudo}}
\sbcend
\sbcend
\end{pilha}
\end{codigo}

\begin{curiosidade}
Existe também a \textbf{variável local}: ``somente para este ator''. Só aquele ator enxerga o valor. Já a variável ``para todos os atores'' é lida por todos os personagens do jogo.
\end{curiosidade}

\section{Hora de praticar}

\begin{exercicio}[Vidas]
Crie a variável \textbf{vidas}, começando com 3. Cada vez que o gatinho tocar na borda, perde 1 vida.
\end{exercicio}

\begin{solucao}
\begin{pilha}
\bk[scEvents]{quando \bandeira\ for clicado}
\bk[scVariables]{defina \menu{vidas} para \entrada{3}}
\sbc[scControl]{sempre}
\sbc[scControl]{se \booleano{tocando em \menu{borda}?} então}
\bk[scVariables]{mude \menu{vidas} em \entrada{$-1$}}
\bk[scControl]{espere \entrada{1} segundos}
\sbcend
\sbcend
\end{pilha}
O \bloco[scControl]{espere} evita que as vidas acabem num instante!
\end{solucao}

\begin{exercicio}[Contador regressivo]
Crie a variável \textbf{tempo} começando em 10. A cada segundo, diminua 1, até chegar a zero.
\end{exercicio}

\begin{solucao}
\begin{pilha}
\bk[scEvents]{quando \bandeira\ for clicado}
\bk[scVariables]{defina \menu{tempo} para \entrada{10}}
\sbc[scControl]{repita até que \booleano{\variavel{tempo} = \entrada{0}}}
\bk[scControl]{espere \entrada{1} segundos}
\bk[scVariables]{mude \menu{tempo} em \entrada{$-1$}}
\sbcend
\bk[scLooks]{diga \entrada{Acabou o tempo!}}
\end{pilha}
\end{solucao}

\begin{exercicio}[Soma da turma]
Pergunte dois números e faça o gatinho dizer a \textbf{soma}. Use duas variáveis, \textbf{a} e \textbf{b}.
\end{exercicio}

\begin{solucao}
\begin{pilha}
\bk[scEvents]{quando \bandeira\ for clicado}
\bk[scSensing]{pergunte \entrada{Primeiro número?} e espere}
\bk[scVariables]{defina \menu{a} para \variavel{resposta}}
\bk[scSensing]{pergunte \entrada{Segundo número?} e espere}
\bk[scVariables]{defina \menu{b} para \variavel{resposta}}
\bk[scLooks]{diga \reporter{\variavel{a} + \variavel{b}} por \entrada{3} segundos}
\end{pilha}
\end{solucao}

\begin{desafio}[Quem chega primeiro a 5?]
Faça um jogo de apertar a barra de espaço: a cada aperto, a variável \textbf{pontos} sobe 1. Quando chegar a 5, o gatinho diz ``Vitória!'' e o jogo para.
\end{desafio}

\begin{solucao}
\begin{pilha}
\bk[scEvents]{quando \bandeira\ for clicado}
\bk[scVariables]{defina \menu{pontos} para \entrada{0}}
\end{pilha}
\begin{pilha}
\bk[scEvents]{quando a tecla \menu{espaço} for pressionada}
\bk[scVariables]{mude \menu{pontos} em \entrada{1}}
\sbc[scControl]{se \booleano{\variavel{pontos} = \entrada{5}} então}
\bk[scLooks]{diga \entrada{Vitória!}}
\bk[scControl]{pare \menu{tudo}}
\sbcend
\end{pilha}
\end{solucao}
